% Required in the document preamble:
% \usepackage{tikz}
% \usetikzlibrary{arrows.meta,positioning,calc}

\chapter{Methodology}
\label{chap:methodology}

This chapter describes the implemented framework for identifying weaknesses in
synthetic tabular health data and using those weaknesses to guide a controlled
retraining intervention. The central design separates two questions. First,
which features contribute most strongly to detectable differences between real
and synthetic records? Second, within those features, which regions of the real
distribution are underrepresented by the generator? Feature-level signals
answer the first question, whereas one-sided region deficits answer the second.
The two are combined only when assigning sampling weights to real training
rows. This distinction is important: SHAP values determine where to focus, but
they do not directly identify which real records should be replicated.

\section{Approach Overview}
\label{sec:approach_overview}

Figure~\ref{fig:methodology_workflow} presents the complete workflow. The real
dataset is divided once into mutually exclusive training, audit, validation,
and test partitions. A baseline generator (A0) is fitted only on the training
partition. Synthetic records from A0 and real records from the audit partition
are then used to fit a post-hoc real-versus-synthetic detector. TreeSHAP,
feature-wise distribution mismatch, and tail or rare-category deficit signals
are derived from this audit comparison. These signals define the feature
priorities for variants A2--A5. Separately, one-sided frequency deficits are
estimated for regions of every feature. Each real training row receives a
weight according to the deficits of the selected-feature regions that it
occupies. Weighted bootstrap samples are appended to the original training
data, and a fresh generator is fitted for each variant. A1 uses the same amount
of additional data but samples uniformly, thereby controlling for the effect
of increasing the training-set size. All generators are evaluated on the same
untouched validation or test partition.

\begin{figure}[htbp]
    \centering
    \begin{tikzpicture}[
        node distance=5mm and 8mm,
        box/.style={draw, rounded corners, align=center, minimum height=9mm,
                    text width=48mm, font=\small},
        branch/.style={box, text width=42mm},
        arrow/.style={-{Latex[length=2mm]}, thick},
        sidearrow/.style={arrow, rounded corners},
        note/.style={font=\scriptsize, align=center, fill=white, inner sep=1pt}
    ]
        \node[box] (data) {Real tabular health data};
        \node[box, below=of data] (split) {Target-stratified partition\\60\% train, 20\% audit,\\10\% validation, 10\% test};
        \node[box, below=of split] (a0) {Fit baseline generator A0\\on the real training partition};
        \node[box, below=of a0] (auditgen) {Generate the A0 synthetic\\audit sample};
        \node[box, below=of auditgen] (signals) {Audit comparison\\random-forest detector and TreeSHAP;\\mismatch and tail-deficit signals};

        \node[branch, below left=of signals.south] (priority) {Feature-level path\\A2--A5 priorities and\\top-$k$ feature selection};
        \node[branch, below right=of signals.south] (regions) {Region-level path\\one-sided deficits between\\real audit and A0};

        \node[box, anchor=north] (weights)
            at ([yshift=-7mm]$(signals.center |- priority.south)$)
            {Training-data intervention\\uniform bootstrap for A1;\\weighted bootstrap for A2--A5};
        \node[box, below=of weights] (retrain) {Fit fresh A1--A5\\generator instances};
        \node[box, below=of retrain] (evaluate) {Common evaluation on validation or test\\fidelity, utility, tails, detectability,\\and privacy proxies};

        \draw[arrow] (data) -- (split);
        \draw[arrow] (split) -- node[note, right] {real train} (a0);
        \draw[arrow] (a0) -- (auditgen);
        \draw[arrow] (auditgen) -- (signals);
        \draw[sidearrow] (split.west) -- ++(-14mm,0) |- node[note, near end, left] {real audit} (signals.west);
        \draw[arrow] (signals) -- (priority);
        \draw[arrow] (signals) -- (regions);
        \draw[arrow] (priority) |- (weights);
        \draw[arrow] (regions) |- (weights);
        \draw[sidearrow] (split.east) -- ++(14mm,0) |- node[note, near end, right] {original real train} (weights.east);
        \draw[arrow] (weights) -- (retrain);
        \draw[arrow] (retrain) -- (evaluate);
        \draw[sidearrow] (a0.east) -- ++(26mm,0) |- node[note, near end, right] {A0 synthetic result} (evaluate.east);
        \draw[sidearrow] (split.west) -- ++(-26mm,0) |- node[note, near end, left] {validation or test} (evaluate.west);
    \end{tikzpicture}
    \caption{Implemented XAI-guided retraining workflow. The audit partition is
    used to construct the intervention; the validation or test partition is
    used only for final comparison of the variants.}
    \label{fig:methodology_workflow}
\end{figure}

\section{Data Partitioning and Experimental Control}
\label{sec:data_partitioning}

Let the complete real dataset be denoted by $D$. A target-stratified split with
fixed seed 42 partitions $D$ into

\begin{equation}
    D = D_{\mathrm{train}} \mathbin{\dot\cup} D_{\mathrm{audit}}
        \mathbin{\dot\cup} D_{\mathrm{val}} \mathbin{\dot\cup} D_{\mathrm{test}},
\end{equation}

with proportions $0.60$, $0.20$, $0.10$, and $0.10$, respectively. The split
indices are persisted and checked for overlap and completeness. The training
partition is the only source of rows used to fit a generator. The audit
partition is reserved for constructing the detector, feature signals, and
region deficits. During model development, variants are evaluated on
$D_{\mathrm{val}}$. Once a configuration is frozen, the same procedure can be
run once against $D_{\mathrm{test}}$. This separation prevents the rows used to
measure final performance from influencing the reweighting intervention.

The same persisted split, generator seed, augmentation size, and evaluation
protocol are used across all variants within a run. Each variant nevertheless
fits a fresh generator instance in an isolated backend directory. Thus, the
comparison changes the composition of the generator's training data while
holding the remaining experimental conditions fixed.

\section{Baseline Synthetic Data Generation}
\label{sec:baseline_generation}

The baseline variant A0 is fitted directly on $D_{\mathrm{train}}$, without
augmentation. It produces two synthetic datasets. The first contains
$|D_{\mathrm{audit}}|$ rows and is used only to construct the intervention. The
second contains $|D_{\mathrm{train}}|$ rows and is used in the common evaluation
protocol. All later feature priorities and region deficits are therefore
defined relative to the same A0 generator.

The implementation exposes CTAB-GAN+, CTGAN, and the DP-CGANS generator
architecture through a common dataframe-in/dataframe-out adapter. This allows
the same split, intervention, and evaluation code to be used for all three
backends. The DP-CGANS backend is explicitly run with
\texttt{private=false}; it is an architectural baseline and must not be
interpreted as a differentially private model or as an end-to-end privacy
guarantee.

Before the retraining ablation, the three generator architectures can be
compared without intervention on the same persisted split. The split seed
remains 42, while generator seeds 42, 43, and 44 provide three independent
fits per architecture. Every completed fit is checkpointed, sampled, and
evaluated using the protocol in Section~\ref{sec:evaluation_protocol}; results
are summarised by their mean and standard deviation. This initial comparison
measures architecture-level differences, whereas the A0--A5 experiment tests
the effect of changing the training-row sampling distribution.

\section{Post-Hoc Detector and Feature-Level Signals}
\label{sec:feature_signals}

\subsection{Real-versus-Synthetic Detector}

A balanced audit probe is formed from equally many rows from
$D_{\mathrm{audit}}$ and the A0 synthetic audit sample. Real records receive
label 1 and synthetic records label 0. The probe is divided, using stratified
sampling, into a 70\% detector-training partition and a 30\% detector holdout.
Numeric variables are median-imputed; categorical variables are
most-frequent-imputed and one-hot encoded. A class-balanced random forest with
300 trees is then trained to estimate

\begin{equation}
    h(\mathbf{x}) = P(y_{\mathrm{real}}=1\mid\mathbf{x}).
\end{equation}

Detector ROC-AUC, average precision, accuracy, and all four real/synthetic
classification outcomes are retained as audit diagnostics.

\subsection{TreeSHAP Importance}
\label{subsec:treeshap_importance}

TreeSHAP explanations are not computed over every detector holdout record.
They are computed for at most 100 \emph{correctly detected synthetic} holdout
records. This scope targets features that expose detectable generator
artifacts. Synthetic records that already fool the detector and all real
records are excluded from the weighting signal.

For a numeric feature $j$, the raw feature importance is

\begin{equation}
    I_j = \frac{1}{|E|}\sum_{i\in E}|\phi_{ij}|,
    \label{eq:raw_shap_importance}
\end{equation}

where $E$ is the explained subset. For a categorical feature, the signed SHAP
contributions of its one-hot columns are first summed within each row; the
absolute value is taken only after this aggregation. The implementation scales
the resulting importances by their maximum, not by their sum:

\begin{equation}
    \widehat{I}_j =
    \begin{cases}
        I_j / \max_k I_k, & \max_k I_k > 0,\\
        0, & \text{otherwise}.
    \end{cases}
    \label{eq:max_normalized_shap}
\end{equation}

If there are no correctly detected synthetic holdout records, the SHAP signal
is explicitly set to zero rather than being replaced by explanations from a
different outcome group.

\subsection{Distribution-Mismatch Signal}
\label{subsec:mismatch_signal}

The mismatch signal measures feature-level disagreement between
$D_{\mathrm{audit}}$ and the A0 synthetic audit sample. For a continuous
feature $j$, the raw mismatch is the first Wasserstein distance divided by the
interquartile range of the real audit feature:

\begin{equation}
    U_j =
    \frac{W_1(X^{\mathrm{audit}}_j,X^{\mathrm{A0}}_j)}
         {\operatorname{IQR}(X^{\mathrm{audit}}_j)}.
    \label{eq:continuous_mismatch}
\end{equation}

If the interquartile range is zero, the standard deviation and then 1 are used
as fallbacks. For a categorical feature, $U_j$ is the Jensen--Shannon distance
between the real and synthetic category distributions. As with SHAP, the
feature values are max-normalized:

\begin{equation}
    \widehat{U}_j = U_j / \max_k U_k,
\end{equation}

with an all-zero result when the maximum is zero.

\subsection{Tail and Rare-Outcome Deficit Signal}
\label{subsec:tail_signal}

For a continuous feature, the raw tail signal is the sum of one-sided deficits
below the real audit 5th percentile and above its 95th percentile:

\begin{equation}
\begin{split}
    T_j ={}& \max\!\left(0,
        P_{\mathrm{audit}}(X_j\le q_{j,.05})
        -P_{\mathrm{A0}}(X_j\le q_{j,.05})\right)\\
        &+\max\!\left(0,
        P_{\mathrm{audit}}(X_j\ge q_{j,.95})
        -P_{\mathrm{A0}}(X_j\ge q_{j,.95})\right).
\end{split}
\label{eq:tail_deficit}
\end{equation}

For a categorical feature, the signal sums one-sided deficits for categories
whose real audit frequency is at most 5\%. For the target feature, the minority
class is included even if its frequency exceeds 5\%. The resulting feature
signals are max-normalized to obtain $\widehat{T}_j$.

The tail signal affects feature selection in A5. It is not an independent
binary multiplier applied to every tail-region row.

\section{Feature Prioritisation and Ablation Variants}
\label{sec:feature_prioritisation}

The three normalized feature signals are combined before any row weights are
computed. For feature $j$, the implemented variant-specific score $C_j^{(a)}$
is

\begin{equation}
    C_j^{(a)} =
    \begin{cases}
        \widehat{U}_j, & a=\mathrm{A2},\\
        \widehat{I}_j, & a=\mathrm{A3},\\
        0.625\widehat{I}_j+0.375\widehat{U}_j, & a=\mathrm{A4},\\
        0.5\widehat{I}_j+0.3\widehat{U}_j+0.2\widehat{T}_j,
            & a=\mathrm{A5}.
    \end{cases}
    \label{eq:variant_feature_score}
\end{equation}

The top $k=5$ features are selected using a stable descending sort. The
selected scores are converted into feature priorities that sum to one:

\begin{equation}
    \rho_j^{(a)} =
    \frac{C_j^{(a)}}{\sum_{k\in\mathcal{J}_a}C_k^{(a)}},
    \qquad j\in\mathcal{J}_a,
    \label{eq:feature_priority}
\end{equation}

where $\mathcal{J}_a$ denotes the selected feature set. Priorities are zero if
the selected-score sum is zero. The code can optionally restrict selection to
one feature per connected correlation group, defined by
$|r|\ge 0.65$, but this option is disabled in the thesis configuration files.

The six controlled variants are therefore:

\begin{table}[htbp]
    \centering
    \caption{Ablation variants implemented in the retraining runner.}
    \label{tab:ablation_variants}
    \begingroup
    \small
    \setlength{\tabcolsep}{4pt}
    \begin{tabular}{p{13mm}p{48mm}p{68mm}}
        \hline
        Variant & Training-data intervention & Purpose \\
        \hline
        A0 & No augmentation & Baseline generator \\
        A1 & Uniform bootstrap augmentation & Controls for added data volume \\
        A2 & Mismatch-prioritised weighting & Distribution mismatch only \\
        A3 & TreeSHAP-prioritised weighting & Detector explanation only \\
        A4 & $0.625$ SHAP $+0.375$ mismatch & Combined XAI and fidelity signal \\
        A5 & $0.5$ SHAP $+0.3$ mismatch $+0.2$ tail & Adds tail/rare deficit signal \\
        \hline
    \end{tabular}
    \endgroup
\end{table}

\section{Underrepresented-Region Modelling}
\label{sec:region_modelling}

Feature priorities identify \emph{which variables} should guide retraining.
Region modelling identifies \emph{which values} of those variables are
underrepresented. The same region definitions are used by A2--A5.

For each continuous feature $j$, bin boundaries are obtained from the unique
training-set quantiles at probabilities $0$, $0.05$, $0.25$, $0.75$, $0.95$,
and $1$, with unbounded exterior intervals. Because the boundaries are learned
from $D_{\mathrm{train}}$, the audit comparison does not redefine the support.
For a categorical feature, each observed category, including an explicit
missing-value category, forms a region.

For region $r\in\mathcal{R}_j$, let $p^{\mathrm{audit}}_{jr}$ and
$p^{\mathrm{A0}}_{jr}$ denote its frequencies in the real and synthetic audit
samples. The implementation uses the absolute one-sided deficit

\begin{equation}
    d_{jr} = \max\left(0,
        p^{\mathrm{audit}}_{jr}-p^{\mathrm{A0}}_{jr}\right).
    \label{eq:region_deficit}
\end{equation}

No relative-frequency alternative is used. To make regions comparable within
each feature, the deficits are divided by the largest deficit of that feature:

\begin{equation}
    g_{jr} =
    \begin{cases}
        d_{jr}/\max_{s\in\mathcal{R}_j}d_{js},
            & \max_s d_{js}>0,\\
        0, & \text{otherwise}.
    \end{cases}
    \label{eq:normalized_region_deficit}
\end{equation}

Consequently, $g_{jr}\in[0,1]$. A region receives a positive score only when
it is less frequent in the A0 sample than in the real audit data; already
matched or overrepresented regions receive zero.

\section{Row Weighting and Bootstrap Augmentation}
\label{sec:row_weighting}

Let $r_j(\mathbf{x}_i)$ denote the region occupied by training row
$\mathbf{x}_i$ for feature $j$. For weighted variant $a\in\{\mathrm{A2},
\mathrm{A3},\mathrm{A4},\mathrm{A5}\}$, the row score is

\begin{equation}
    s_i^{(a)} = \sum_{j\in\mathcal{J}_a}
        \rho_j^{(a)}g_{j,r_j(\mathbf{x}_i)}.
    \label{eq:implemented_row_score}
\end{equation}

Because the non-negative feature priorities sum to one and the region scores
lie in $[0,1]$, $s_i^{(a)}$ also lies in $[0,1]$. No additional division by
the number of selected features or normalization by the largest observed row
score is applied. The final weight is

\begin{equation}
    w_i^{(a)} = \min\left(1+\alpha s_i^{(a)},w_{\max}\right),
    \label{eq:implemented_row_weight}
\end{equation}

with $\alpha=4$ and $w_{\max}=3$ in the thesis configurations. Sampling
probabilities are

\begin{equation}
    \pi_i^{(a)} = \frac{w_i^{(a)}}{\sum_k w_k^{(a)}}.
    \label{eq:implemented_sampling_probability}
\end{equation}

For every augmented variant, $\lfloor\gamma|D_{\mathrm{train}}|\rfloor$
rows are sampled with replacement and appended to the complete original
training partition. The configured augmentation ratio is $\gamma=0.6$.
Variants A2--A5 sample according to Equation~\ref{eq:implemented_sampling_probability};
A1 samples the same number of rows uniformly. The intervention therefore
never removes a real training row. A weight of 1.3 represents a relative
sampling propensity, not 1.3 deterministic copies of a record.

\section{Generator Retraining}
\label{sec:generator_retraining}

Each A1--A5 dataset is used to fit a fresh instance of the selected generator.
The fitted model then produces exactly $|D_{\mathrm{train}}|$ synthetic rows
for evaluation, irrespective of the augmented training-set size. This keeps
the number of evaluated synthetic records constant across variants.

CTAB-GAN+ uses its mixed-type data transformer, convolutional generator and
discriminator, conditional-vector loss, information loss, and auxiliary
task-classifier loss. For MIMIC-IV it is trained for 150 epochs with batch size
512, latent dimension 100, and 64 channels. For WiDS it is trained for 200
epochs with batch size 1024, latent dimension 128, and 96 channels. CTGAN uses
two 256-unit hidden layers in both generator and discriminator, a 128-dimensional
embedding, one discriminator step, and packing degree 10. The DP-CGANS
architecture uses three 128-unit hidden layers, ten discriminator steps, and
packing degree 10, always in non-private mode. CTGAN and DP-CGANS use batch
sizes 500 for MIMIC-IV and 1000 for WiDS, with 150 and 200 epochs,
respectively.

All adapters infer numeric measurement precision and finite training support
from the rows fitted by that variant. Generated numeric values are restored to
the fitted dataframe schema and rounded to the inferred measurement precision;
raw pre-restoration samples and post-processing diagnostics are saved. For
DP-CGANS ablations, A0 fits the transformer once from the original real
training partition. A hash-validated frozen copy is reused by A1--A5, while
each variant still transforms its own augmented rows, rebuilds its sampler,
and trains a fresh GAN. This prevents transformer changes from becoming an
uncontrolled source of variation.

\section{Evaluation Protocol}
\label{sec:evaluation_protocol}

Every variant is compared with the same real validation or test partition.
Evaluation has five complementary components.

\begin{enumerate}
    \item \textbf{Fidelity.} Continuous features are assessed using
    Wasserstein distance, IQR-scaled Wasserstein distance, the
    Kolmogorov--Smirnov statistic, and correlation-matrix distance.
    Categorical features are assessed using Jensen--Shannon distance.

    \item \textbf{Tail and rare-event coverage.} Continuous-feature lower and
    upper quantile mass, scaled quantile error, and CDF-tail divergence are
    measured. Rare-category frequency error and minority-outcome frequency
    error are also retained.

    \item \textbf{Downstream utility.} A class-balanced random forest is
    trained on synthetic data and evaluated on untouched real data. The fixed
    argmax decision rule is used without threshold tuning. Mortality is tested
    both with the real training prevalence and with a 50/50 training balance.
    ROC-AUC, PR-AUC, accuracy, balanced accuracy, macro precision, macro
    recall, macro F1, and positive-class metrics are reported. Additional
    additive and replacement curves measure utility for controlled mixtures of
    real and synthetic training rows.

    \item \textbf{Detectability.} A new real-versus-synthetic random forest is
    fitted between the evaluation partition and each variant's synthetic data.
    Lower separability, interpreted relative to the chance target of 0.5,
    indicates improved realism.

    \item \textbf{Privacy proxies.} Exact matches with the real training data
    and nearest-neighbour distances from synthetic and held-out real records to
    the real training set are reported. These are memorisation-risk indicators,
    not formal privacy guarantees.
\end{enumerate}

The primary controlled reference is A0. Results therefore report A1--A5 minus
A0 for fidelity, tail, utility, detector, and privacy-proxy measures. Utility
is additionally compared with a real-only model trained on
$D_{\mathrm{train}}$ and evaluated on the identical frozen split.

\section{Internal-Discriminator Explanation Experiment}
\label{sec:internal_discriminator_experiment}

The internal-discriminator explanation experiment is separate from the
XAI-guided reweighting framework described in
Sections~\ref{sec:approach_overview}--\ref{sec:evaluation_protocol}. Its purpose
is to investigate how the CTAB-GAN+ discriminator's decisions and feature
attributions develop during adversarial training. Its outputs are not used to
select features, define underrepresented regions, calculate row weights, or
train any A0--A5 variant. The reweighting intervention relies exclusively on
the post-hoc random-forest detector and its TreeSHAP explanations.

\subsection{Snapshot Collection and Common Probe}

CTAB-GAN+ discriminator snapshots are captured every 25 epochs in the
MIMIC-IV configuration. Snapshot collection is disabled in the WiDS
configuration. Each historical discriminator is evaluated against the same
fixed, balanced probe constructed from held-out real audit records and records
from the final generator of the analysed variant. Holding the probe constant
isolates changes in the discriminator: the resulting trajectory shows how
historical critics respond to one common final-generator sample and is not an
epoch-matched generator--discriminator comparison.

The probe is divided into disjoint calibration and holdout partitions. A
real-versus-synthetic decision threshold is selected only on the calibration
partition by maximising Youden's $J$ statistic and is subsequently applied to
the holdout partition. Threshold-independent ROC-AUC and average precision are
reported alongside threshold-dependent classification outcomes.

\subsection{GradientSHAP Analysis}

GradientSHAP explanations are computed separately for four holdout outcome
groups: correctly classified real records, real records classified as
synthetic, correctly classified synthetic records, and synthetic records
classified as real. Valid CTAB-GAN+ conditional vectors are sampled and
marginalised over eight draws instead of evaluating the discriminator with an
out-of-distribution all-zero condition. The experiment uses up to 50
calibration records as the explanation background and up to 100 holdout records
per explanation scope.

The principal cross-model comparison uses the correctly classified synthetic
group, which matches the scope of the post-hoc detector's TreeSHAP analysis.
For every snapshot, the experiment stores discriminator metrics, row-level
outcomes, signed and absolute feature attributions, and feature-attribution
trajectories. At the final snapshot, the GradientSHAP ranking is compared with
the post-hoc random-forest TreeSHAP ranking. This comparison assesses agreement
between internal and external explanations; it does not feed information back
into generator retraining.

\section{Implementation Details and Reproducibility}
\label{sec:implementation_details}

The framework is implemented in Python. Data manipulation uses pandas and
NumPy; statistical distances use SciPy; preprocessing, random forests,
stratified splitting, and nearest-neighbour analysis use scikit-learn;
explanations use SHAP's \texttt{TreeExplainer}; and the GAN implementations use
PyTorch. The standalone generator dependencies are CTGAN 0.12.1 and DP-CGANS
0.2.0. Experiments are launched through the
\texttt{xai\_reweighting.run\_ablation} command-line runner or its Jupyter
orchestration notebook.

Random seeds are propagated to Python, NumPy, and PyTorch. Deterministic
execution is enabled, TF32 is disabled in the thesis configurations, and the
selected device is recorded. Configuration, source-data, and code hashes form
a run fingerprint. Split indices, feature scores, region definitions, row
weights, augmentation counts, synthetic datasets, model checkpoints, training
histories, convergence diagnostics, and metric files are written to the run
directory. Writes are atomic, and resuming a run is permitted only when the
stored fingerprint matches the current configuration, data, and code.

CTAB-GAN+'s continuous-column mixture preprocessing retries a non-converged fit
and stops if convergence still fails. Non-finite losses or model parameters also stop
the experiment. These checks prevent failed training runs from silently
entering the variant comparison.
